% element: 10-s018:e03
\BeleskeID{10-s018}
Polazimo od aktivne karakteristike
\[V_O=V_{CC}-\beta_F\frac{R_C}{R_B}(V_I-V_{BE}).\]
U pragu prebacivanja ulaz i izlaz imaju isti napon, pa zamenjujemo oba sa $V_M$ i sakupljamo članove uz $V_M$. Konačni izraz može se zapisati u još dva ekvivalentna oblika:
\[\begin{aligned}
V_M&=V_{BE}\frac{\frac{V_{CC}}{V_{BE}}+\beta_F\frac{R_C}{R_B}}{1+\beta_F\frac{R_C}{R_B}}\\
&=V_{BE}\left(\frac{\frac{V_{CC}}{V_{BE}}-1}{1+\beta_F\frac{R_C}{R_B}}+1\right).
\end{aligned}\]
Drugi oblik pokazuje da se pri velikom pojačanju prag približava $V_{BE}$. Jednostruke margine određuju se u odnosu na taj prag, zbog čega je margina logičke nule kod RTL kola mala.
